\section{Cofibres, reversible excision, and integral homology}

\begin{claim}{B.1}{Known}{The baseline used in this part}
For the modern fundamental localizers fixed on the first page, $\Winfty\subseteq\W$, and $\W$ is strongly saturated and closed under retracts in $\Cat^{[1]}$. Neither assertion requires accessibility. Consequently a specified zigzag of weak homotopy equivalences of categories preserves $\W$-asphericity, and a commutative square whose vertical arrows are in $\Winfty$ transfers membership of its horizontal arrows in $\W$ in both directions.
\end{claim}
The first assertion is \citet[Theorem 2.2.11, p.~136]{RootCisinski2004}; the second is \citet[Proposition 4.2.4, p.~183]{RootCisinski2006}. For the last assertion, write the equality of the two composites in the square and apply two out of three. For a zigzag, apply this argument to each of its arrows. We do not identify this modern definition with the weaker printed XVI Loc (1)--(3); Part A records the discrepancy.

\begin{claim}{B.2}{Proved}{Forward gluing is automatic under relative Theorem A}
Let $J$ be a small category, $F,G:J\to\Cat$ covariant functors, and $\eta:F\Rightarrow G$ a natural transformation with $\eta_j\in\W$ for every $j$. The functor $\int_J\eta:\int_JF\to\int_JG$ belongs to $\W$. In particular every modern fundamental localizer satisfies Loc (6 a), and the augmentation from the Grothendieck construction of a sieve-cover span to its union is in $\Winfty$.
\end{claim}
\begin{proof}
An object of $\int_JF$ is $(j,x)$, $x\in F(j)$; an arrow $(j,x)\to(k,y)$ is $(\alpha,v)$ with $\alpha:j\to k$ and $v:F(\alpha)x\to y$. Composition is
\[
 (\beta,w)(\alpha,v)=(\beta\alpha,w\,F(\beta)(v)).
\]
Over an object $s\in J$, evaluation
\[
 E_s:(\int_JF)/s\longrightarrow F(s),\qquad
 (j,x,\alpha:j\to s)\longmapsto F(\alpha)x
\]
is left adjoint to $I_s(y)=(s,y,\id_s)$, with identity counit. The unit has component $(\alpha,\id)$. These adjunctions give specified nerve homotopy equivalences, naturally in $F$. The square with horizontal maps induced by $\eta$ and vertical maps $E_s$ commutes. Thus $(\int_J\eta)/s\in\W$ by two out of three, and relative Theorem A proves the assertion. This is also the explicit argument behind \citet[Proposition 1.2.4]{RootCisinski2004}.

For a span $U\leftarrow K\to V$, its Grothendieck construction has three layers $K,U,V$. The subcategories $K\cup U$ and $K\cup V$ are full sieves, each retracting by the displayed adjunction onto its outer layer. Their intersection is $K$. Their nerves include $NK$ by simplicial monomorphisms and cover the nerve of the entire construction, so that union is a homotopy pushout. The retractions give a map from the span $(K\cup U)\leftarrow K\to(K\cup V)$ to $U\leftarrow K\to V$, identity on the middle category and a weak equivalence on each arm. Homotopy-pushout invariance therefore supplies the specified weak-equivalence comparison with the double mapping cylinder. No equality of the nerve with a simplicial replacement is asserted. For arbitrary indexing categories, the general Grothendieck-construction comparison used below is Thomason's homotopy-colimit theorem, in the natural form proved by \citet[Theorem 26.8 and Corollary 26.9]{RootChacholskiScherer2001}.

Now suppose $K=U\cap V$ and $X=U\cup V$ is a cover by full sieves. A chain in $X$ lies in a piece containing its last object. Hence $NX=NU\cup_{NK}NV$ literally, and both attaching maps are simplicial monomorphisms. Realization gives inclusions of CW subcomplexes; the augmentation from the double mapping cylinder to this union is a homotopy equivalence. Thus the specified functor
\[
 q:\int(U\leftarrow K\to V)\longrightarrow X
\]
which forgets each layer tag belongs to $\Winfty$. A map between two cover squares induces a natural transformation of their spans. If its three input maps are in $\W$, the first paragraph and the two augmentation functors show that its union map is in $\W$. This is exactly Loc (6 a). Empty pieces cause no change: the same statements hold for the empty simplicial subcomplex and the disjoint union.
\end{proof}

\begin{claim}{B.3}{Proved}{A functorial categorical homotopy cofibre}
For every functor $g:A\to B$ define $M_g$ on disjoint layers $A$ and $B$, with their original internal arrows and
\[
 \Hom_{M_g}(a,b)=\Hom_B(g(a),b),\qquad
 \Hom_{M_g}(b,a)=\varnothing.
\]
Composition of crossing arrows is induced by $g$ and composition in $B$. Adjoin an object $t$, with $\Hom(a,t)=\{*\}$ for every $a\in A$ and no arrows between $t$ and the $B$-layer. Call the resulting category $K(g)$. There is a functorial homotopy-pushout square of spaces
\[
 \begin{tikzcd}
 BA\arrow[r,"Bg"]\arrow[d]&BB\arrow[d]\\
 *\arrow[r]&BK(g),
 \end{tikzcd}
\]
where $t$ supplies the basepoint. Thus $K(g)$ represents the unpointed homotopy cofibre, based at its cone point. A commutative square of functors induces a specified functor of these categories. If both vertical functors are in $\W$, the induced functor on $K(-)$ is in $\W$.
\end{claim}
\begin{proof}
The formulas define a category: the only mixed associativity equations are the functoriality of $g$ or associativity in $B$. The retraction $r_g:M_g\to B$, equal to $g$ on the first layer and the identity on the second, is left adjoint to the inclusion $s_g:B\hookrightarrow M_g$. Its unit is the crossing arrow $a\to g(a)$ corresponding to $\id_{g(a)}$, and its counit is the identity. In particular $r_g\in\Winfty$ and $r_gi_g=g$, where $i_g:A\hookrightarrow M_g$ is the other inclusion.

The full subcategories $M_g$ and $A^{\triangleright}=A\cup\{t\}$ cover $K(g)$ by sieves and meet in $A$. The argument of \claimref{B.2} identifies their union with the homotopy pushout of the specified span. The map of spans $(r_g,\id_A,A^{\triangleright}\to e)$ is objectwise a weak homotopy equivalence. It identifies that pushout with the displayed square, including the prescribed map $Bg$, rather than merely its source and target homotopy types.

For a commutative square $v g=g'u$ with $u:A\to A'$ and $v:B\to B'$, the induced functor sends the layers by $u,v$, sends $t$ to $t'$, and sends a crossing arrow $g(a)\to b$ by $v$. This is well-defined because $vg=g'u$. It is the map of the Grothendieck constructions of the spans $(B\leftarrow A\to e)$ and $(B'\leftarrow A'\to e)$. Apply \claimref{B.2}. In particular no condition on components, endomorphisms, or idempotents was used.

If $A=\varnothing$, then $K(g)=B\amalg e$. If also $B=\varnothing$, then $K(g)=e$. There is no functor from a nonempty $A$ to an empty $B$. These are precisely the homotopy-pushout conventions needed in degrees zero and minus one; the empty source is not silently replaced by a chosen basepoint.
\end{proof}

\begin{claim}{B.4}{Proved}{The manuscript cylinder and the exact cofibration hypothesis}
Let $\Cof_\W$ mean the functors $i:A\to B$ for which $B\amalg_A(-):A\backslash\Cat\to B\backslash\Cat$ preserves every $\W$-map. Every classical Dwyer map is in $\Cof_\W$. Every functor $g:A\to B$ has a functorial factorization $g=p_gj_g$, with $j_g$ a Dwyer map and $p_g\in\Winfty$. The ordinary pushout $e\amalg_A\operatorname{Cyl}(g)$ is a categorical homotopy cofibre of $g$.
\end{claim}
\begin{proof}
A Dwyer map is a full sieve inclusion $A\subseteq B$ such that its inclusion into the cosieve $V$ generated by $A$ has a right adjoint $r:V\to A$ restricting to the identity. The pushout comparison
\[
 \int(B\leftarrow A\to C)\longrightarrow B\amalg_A C
\]
is in $\Winfty$: this is the Dwyer pushout theorem, \citet[Proposition 4.3]{RootThomason1980}. Its hypotheses concern a Dwyer map, not an arbitrary sieve. Given $C\to C'$ in $\W$ under $A$, apply \claimref{B.2} to the map of spans and use these two comparison functors. This proves the full undercategory definition of $\Cof_\W$, including preliminary cobase changes.

Take the exact interval of XII ms.~228,
\[
 I=\Phi_2=(0\to1\leftarrow2\to3\leftarrow4).
\]
The endpoint inclusion $\{0,4\}\hookrightarrow I$ is a Dwyer map: its generated cosieve has the disjoint components $0\to1$ and $4\to3$, each retracting onto its endpoint by a right adjoint. Taking its product with $A$ preserves the explicit adjunction. Put
\[
 \operatorname{Cyl}(g)=B\amalg_{A\times\{0\}}(A\times I).
\]
The functor $p_g:\operatorname{Cyl}(g)\to B$ is induced by $\id_B$ and $g\operatorname{pr}_A$, and $j_g$ is the inclusion of the free endpoint $A\times\{4\}$. Dwyer maps are closed under cobase change and composition, as follows directly from their right-adjoint description (or \citealp[Proposition 4.3 and Lemma 5.3(1)]{RootThomason1980}). Applying this to the two-endpoint inclusion, then the second summand $A\hookrightarrow B\amalg A$, proves that $j_g$ is Dwyer. The map $B\to\operatorname{Cyl}(g)$ is the cobase change of the Dwyer weak equivalence $A\times\{0\}\to A\times I$. The pushout theorem makes it a weak homotopy equivalence. Since its composite with $p_g$ is $\id_B$, $p_g\in\Winfty$.

The same theorem identifies the indicated ordinary pushout along $j_g$ with the homotopy pushout of $B\leftarrow A\to e$. A specified comparison zigzag with $K(g)$ is obtained from the map of span Grothendieck constructions induced by $p_g$ and from their augmentation functors. Every arrow in that zigzag is in $\Winfty$. This justifies using either categorical model in tests of $\W$-asphericity.

Nothing here uses a pushout-product. The product $A\times\{0,4\}\to A\times I$, its cobase change, and its composite with a summand inclusion are the three operations actually used by the manuscript. Loc (6 a) supplies the forward gluing in the manuscript's weaker framework; for our modern localizer it is already \claimref{B.2}. No converse excision is needed to construct a cofibre or to obtain enough $\W$-cofibrations.
\end{proof}

\begin{claim}{B.5}{Proved}{The exact reversible-cofibre criterion: Lead 1}
For a modern fundamental localizer $\W$, the following conditions are equivalent: Loc (6 b); its single-square reverse implication Loc ($6'$ b); the biconditional Loc (6 C) for every $\W$-homotopy-cocartesian square; and
\begin{equation}\label{eq:cofibercriterion}
 g\in\W\quad\Longleftrightarrow\quad K(g)\longrightarrow e\in\W
 \qquad\text{for every small-category functor }g.
\end{equation}
Under these conditions Loc (6 a,b,c), W(7), and W(7 bis) all hold. Accessibility is unnecessary. The implication from left to right in \eqref{eq:cofibercriterion} holds for every fundamental localizer.
\end{claim}
\begin{proof}
If $g\in\W$, the map of spans from $(B\leftarrow A\to e)$ to $(B\leftarrow B\to e)$ induces $K(g)\to K(\id_B)$ in $\W$ by \claimref{B.2}. The latter target is weakly contractible; this proves the unconditional direction.

Suppose Loc ($6'$ b) holds and $K(g)$ is $\W$-aspherical. In the actual sieve-cover square
\[
 \begin{tikzcd}
 A\arrow[r,"i_g"]\arrow[d]&M_g\arrow[d]\\
 A^{\triangleright}\arrow[r,"j"]&K(g),
 \end{tikzcd}
\]
the functor $j$ is in $\W$, since its source has a terminal object and its target is $\W$-aspherical. Reverse excision gives $i_g\in\W$. Then $g=r_gi_g\in\W$. This proof uses no assertion that $i_g$ is itself a $\W$-cofibration; it need not be one.

Conversely suppose \eqref{eq:cofibercriterion}. For a homotopy-pushout square of spaces, the map between the cofibre of its top edge and the cofibre of its bottom edge is a weak equivalence. Indeed both are the colimit of the same expanded diagram after adjoining a cone point, with the two orders of taking pushouts interchanged. In the categorical models just constructed these are the augmentation maps from the iterated span Grothendieck construction: objects are triples consisting of the two span indices and an object in the corresponding category, and the canonical interchange of the indices preserves the arrow-composition formula in \claimref{B.2}. Thus the comparison is attached to the square, with no choice of an abstract isomorphism of homotopy types.

For a $\W$-homotopy-cocartesian square $(A\to B,A\to C,B\to D,C\to D)$, factor its comparison as $P=\int(B\leftarrow A\to C)\to D\in\W$. The preceding weak comparison and \claimref{B.3} give a zigzag of $\W$-maps between $K(A\to B)$ and $K(C\to D)$. The latter step uses the commutative square $C\to P\to D$ with identity on $C$. Their asphericities agree, and \eqref{eq:cofibercriterion} gives Loc (6 C). Sieve-cover squares are such squares by \claimref{B.2}, giving Loc ($6'$ a,b).

We spell out the cube implication, since a single-edge statement must not be substituted for it. For a morphism $f_i:X_i\to Y_i$ of two homotopy-pushout squares, put $C_i=\cofib((BX_i)_+\to(BY_i)_+)$ in pointed spaces. Its underlying space is the homotopy cofibre in \claimref{B.3}. Adding a disjoint basepoint and taking a pointed cofibre are colimit functors. Consequently the four $C_i$ form a homotopy-pushout square of pointed spaces. The forgetful functor to spaces creates these pushouts, whose indexing category is weakly contractible, so their underlying square is also homotopy cocartesian. The comparison is the same iterated-colimit comparison just specified. If $f_1,f_2,f_3\in\W$, then $C_1,C_2,C_3$ are $\W$-aspherical. Applying forward gluing to $C_1\leftarrow C_0\to C_2$ and its map to $*\leftarrow C_0\to *$ makes
\[
 C_3\longrightarrow\Sigma C_0
\]
a $\W$-equivalence. Thus $\Sigma C_0$ is $\W$-aspherical. Applying \eqref{eq:cofibercriterion} to $C_0\to *$ gives asphericity of $C_0$, and applying it to $f_0$ gives $f_0\in\W$. Here statements about spaces are interpreted by their simplex categories; the natural last-vertex comparisons of \citet[Proposition 2.2.3 and Corollary 2.2.4]{RootCisinski2004} transfer the specified comparisons to $\Cat$. This proves Loc (6 b). The implication Loc (6 b)$\Rightarrow$Loc ($6'$ b) is the degenerate cube in Part A.

Finally, if $f_0,f_3\in\W$, then $C_0,C_3$ are $\W$-aspherical. The basepoint section $*\to C_0$ is in $\W$, since its composite with $C_0\to *$ is the identity. The induced map of spans $(C_1\leftarrow *\to C_2)\to(C_1\leftarrow C_0\to C_2)$ and forward gluing give the specified $\W$-map $C_1\vee C_2\to C_3$, so that wedge is $\W$-aspherical. Each summand is a retract, by collapsing the other summand to the specified basepoint. Retract closure from \claimref{B.1} makes each $C_i$, $i=1,2$, aspherical. The criterion gives Loc (6 c). All these cofibre spaces have a specified basepoint, including when an original category is empty. The argument never chooses a point in an empty source.
\end{proof}

\begin{claim}{B.6}{Proved}{The two-cofibre argument: Lead 2}
Every modern fundamental localizer satisfying Loc (6 b), or only Loc ($6'$ b), contains $\WZ$. The conclusion does not require adding Loc (6 a) as an independent axiom.
\end{claim}
\begin{proof}
Let $g:A\to B$ induce an isomorphism on integral homology in every degree, including degree zero. If $A=\varnothing$, then $H_0(BB;\mathbb Z)=0$, so $B=\varnothing$ and $g$ is the identity. Otherwise $BA$ and $BB$ are nonempty and $\pi_0(Bg)$ is a bijection: $H_0(Bg)$ is the free abelian-group map on the actual component function, whose matrix has exactly one $1$ in each column; invertibility forces that function to be a bijection, also for infinite sets.

The homotopy cofibre $C=BK(g)$ is connected. More generally its components consist of one component containing the cone point and all components of $BB$ met by $BA$, together with each unmet component of $BB$. The reduced cofiber exact sequence, or the augmented cellular mapping-cone complex, now gives $\widetilde H_*(C;\mathbb Z)=0$. Thus its unreduced suspension $\Sigma C$ is simply connected by van Kampen and integrally acyclic by the suspension homology sequence. Both are CW spaces. If $\Sigma C$ had a first nonzero homotopy group, the Hurewicz theorem would identify it with a nonzero homology group. Therefore it is weakly contractible; Whitehead's theorem also makes it contractible. These are the standard consequences in \citet[Theorem 4.32 and Corollary 4.33]{RootHatcher2002}.

The categorical cofibre $K(K(g)\to e)$ has precisely that weak homotopy type by \claimref{B.3}, so its map to $e$ belongs to $\Winfty\subseteq\W$. The criterion \claimref{B.5}, applied first to $K(g)\to e$ and then to $g$, gives $g\in\W$. The argument permits arbitrarily many source components; only their bijective correspondence with target components was used. The independent input is reverse excision. Forward gluing and the cofibre construction already follow from relative Theorem A and minimality, as proved above.
\end{proof}

\begin{claim}{B.7}{Proved}{The finite acyclic category, with a complete certificate}
There is a finite category $A$ with $35$ objects and $159$ arrows, all identities included, for which $BA$ is connected and integrally acyclic and $\pi_1(BA)\ne1$.
\end{claim}
\begin{proof}
Use the words $w_0=aaaaaBBB$ and $w_1=bbbBABA$, where capitals mean inverses; their lengths are $8$ and $7$. Define a semisimplicial set with vertices $v,c_0,c_1$, loop edges $a,b$ at $v$, radial edges $r_{k,i}:c_k\to v$ for $i\in\mathbb Z/|w_k|$, and triangles $t_{k,i}$. For a positive letter $\ell$ in position $i$ set
\[
 (d_0t_{k,i},d_1t_{k,i},d_2t_{k,i})=(\ell,r_{k,i+1},r_{k,i});
\]
for $\ell^{-1}$ set this triple to $(\ell,r_{k,i},r_{k,i+1})$. The endpoints of a radial edge are $(v,c_k)$, and those of a loop edge are $(v,v)$. The semisimplicial identities follow by taking endpoints; there are no higher simplices.

An object of $A$ is one of these simplices. An arrow $\sigma\to\tau$ is an order-preserving injection $\theta:[\dim\sigma]\hookrightarrow[\dim\tau]$ with $\theta^*\tau=\sigma$. Composition is composition of injections. This specifies all arrows and composites, including distinct occurrences of the same face. The nerve is the barycentric subdivision of this semisimplicial realization: a nondegenerate nerve simplex is a flag of proper face injections, whose barycentres give the corresponding subdivided simplex. These maps agree with every face identification.

The realization is the cone-triangulation of the presentation complex
\[
 \langle a,b\mid a^5b^{-3},\ b^3(ab)^{-2}\rangle.
\]
Its cellular boundary in degrees two and one is
\[
 d_2=\begin{pmatrix}5&-2\\-3&1\end{pmatrix},\qquad \det d_2=-1,
\]
and $d_1=0$. This proves integral acyclicity. To certify nontrivial fundamental group, map $a$ to $(1\ 4\ 3\ 2\ 5)$ and $b$ to $(1\ 2\ 3)$ in $A_5$, composing right to left. The relators map to the identity: $a^5=b^3=1$ and $ab=(1\ 5)(3\ 4)$. The image of $a$ is nonidentity. This is the finite model of the acyclic complex in \citet[Example 2.38, p.~142]{RootHatcher2002}; the categorical specification above rechecks the supplied TOTASPH example.

The executable \texttt{code/finite\_examples.py} constructs every arrow and composite and checks every identity and associativity equation. It finds normalized nerve-chain ranks $(35,124,90)$ and boundary ranks $(34,90)$. Every nonzero Smith factor is $1$. Thus the images of the boundaries are primitive, and the integral homology is $\mathbb Z$ in degree zero and zero otherwise. The script checks the two permutations separately. The full arrow and composition tables, integer invariants, and output are included; the computer check is not being used to infer contractibility from acyclicity.
\end{proof}

\begin{claim}{B.8}{Refuted}{The reverse axiom for $\Winfty$, including Dwyer cobase changes}
The assertions that $\Winfty$ satisfies W(7), W(7 bis), Loc (6 b), or Loc (6 c) are false. XII Theorem 2 at ms.~78 and XVI Theorem 5.4 insofar as it includes Loc (6) are false. The reverse edge implications can fail even for a pushout along a classical Dwyer map.
\end{claim}
\begin{proof}
Take $A$ of \claimref{B.7}. Adjoin two terminal cone points $u,t$, with one arrow from each object of $A$ to each cone point and no arrows between the cone points. Call this category $X$, and put $U=A\cup\{u\}$ and $V=A\cup\{t\}$. The two pieces are full sieves, cover $X$, and meet in $A$. Their nerves are cones, while $BX$ is the unreduced suspension of $BA$. It is simply connected and acyclic, hence contractible by the argument in \claimref{B.6}. Thus $U,V,X$ are weakly contractible and $A$ is not. The map of this cover square to the constant point square violates the cube reverse implication. The edge $U\to X$ is in $\Winfty$ while $A\to V$ is not, violating W(7 bis). The one-point complement $t$ has $U/t\cong A$, refuting the unqualified one-point converse. The program verifies the $37$ objects and $231$ arrows, both sieve conditions, and integral homology of $X$.

For the Dwyer statement use $g:A\to e$ and the factorization of \claimref{B.4}. The sieve $j_g:A\to\operatorname{Cyl}(g)$ is Dwyer, its target is weakly contractible, and $j_g\notin\Winfty$. The pushout square
\[
 \begin{tikzcd}
 A\arrow[r,"j_g"]\arrow[d]&\operatorname{Cyl}(g)\arrow[d]\\
 e\arrow[r]&e\amalg_A\operatorname{Cyl}(g)
 \end{tikzcd}
\]
has weakly contractible lower-right vertex, because it represents $\Sigma BA$. Both right and bottom arrows are weak equivalences, whereas both left and top arrows are not. This refutes both biconditionals printed in XII ms.~228, footnote 124, when their additional conservative axioms are omitted.

For completeness the finite model used in the code has objects $y$ and $(a,i)$ for $i=1,2,3,4$. Within and between the layers the arrows are those of $A\times\{1\leftarrow2\to3\leftarrow4\}$; from $y$ to each $(a,1)$ there is a unique arrow. There are no other arrows. The pushout replaces the entire fourth layer by a new object $t$, with a unique arrow $t\to(a,3)$. Composition inside product layers uses composition in $A$; composition after either new arrow is the unique new arrow. There are $141$ objects and $1149$ arrows before this quotient, and $107$ objects and $867$ arrows after it. The code checks the quotient functor and all category axioms. The generated cosieve of the fourth layer consists of the third and fourth layers; forgetting their index is the displayed right adjoint, whose Hom-set bijections are checked. This is an explicit finite counterexample, not an appeal to an unspecified cofibrant replacement.
\end{proof}

\begin{claim}{B.9}{Proved}{The two variances of the one-point collage}
For a covariant functor $H:A\to\mathbf{Set}$, let $A_H$ be obtained by adjoining $x$ with $\Hom(x,a)=H(a)$, no arrows $a\to x$, and $\End(x)=\{\id\}$. If $E=\El_A(H)$ and $\pi:E\to A$, there is a literal isomorphism of simplicial sets
\[
 NA_H\cong NA\amalg_{NE}(\Delta^0*NE).
\]
Thus $BA_H$ is the homotopy cofibre of $B\pi$. Dually, for $H:A^{\op}\to\mathbf{Set}$, adjoining a last object $t$ with $\Hom(a,t)=H(a)$ gives the right-cone formula with $NE*\Delta^0$. Under the equivalent axioms in \claimref{B.5}, a full sieve $U\subset X$ with one complementary object $t$ and $\End_X(t)=\{\id\}$ satisfies
\[
 U\to X\in\W\quad\Longleftrightarrow\quad U/t\to e\in\W.
\]
\end{claim}
\begin{proof}
A nondegenerate simplex not in $NA$ begins with $x$, followed by an element $h\in H(a_0)$ and a string $a_0\to\cdots\to a_n$. These are exactly the data of a simplex of $NE$; successive elements are determined by $H$. The first face forgets $h$ and is $N\pi$, while all other faces and degeneracies agree with the cone. This proves the literal simplicial isomorphism, including repeated identities. Realization of $NE\hookrightarrow\Delta^0*NE$ is a CW cofibration, so this ordinary pushout is a homotopy pushout. Passing to opposites gives the right-cone formula, with its variance changed as stated.

For the sieve use $H(u)=\Hom_X(u,t)$, whose category of elements is exactly $U/t$. The square $(U/t\to U,U/t\to e,U\to X,e\to X)$ is a homotopy-pushout square by the formula. Loc (6 C) identifies membership of its top and bottom parallel pair, or of its other parallel pair. The latter says $U\to X\in\W$ iff $U/t\to e\in\W$. This also includes $U/t=\varnothing$, when the attached point is isolated. The one-point statements from the supplied TOTASPH answer are therefore verified with their full stated axiom, and their failure without it is \claimref{B.8}.
\end{proof}

\begin{claim}{B.10}{Proved}{An explicit nontrivial acyclic group}
Every small group $H$ embeds in a small acyclic group $M(H)$ of cardinality at most $\max(|H|,\aleph_0)$; for infinite $H$ equality holds. In particular $M(\mathbb Z)$ is countable, nontrivial, and acyclic. Its one-object category has no nonidentity idempotents.
\end{claim}
\begin{proof}
Form $m(H)$ from $H\times H$ by two HNN extensions with stable letters $d,s$ and relations
\[
 d^{-1}(h,1)d=(h,h),\qquad s^{-1}(h,1)s=(1,h)\quad(h\in H).
\]
The associated subgroups in both extensions are isomorphic to $H$. The normal-form theorem for HNN extensions embeds the base group; equivalently a reduced word with stable letters cannot equal a nontrivial base-group word unless an adjacent stable-letter pair pinches through its specified associated subgroup. Thus these are injections. The resulting group is generated by the first copy of $H$, $d$, and $s$; it has the relations $h^d=hh^s$ and $[h',h^s]=1$. Iterate these injections and take the union $M(H)=\bigcup_{r\geq0}m^r(H)$. Cardinality and the embedding assertion follow from words of finite length in the indicated generators.

We include the homology argument rather than relying on the earlier AI answer. Let $f:L\to H$ kill positive homology over a field $k$ below degree $n$. The homomorphism $q:H\times H\to m(H)$ given by $q(h,h')=h(h')^s$ is defined by the commutation relations. Its restrictions to the two factors and to the diagonal are conjugate to the inclusion $j:H\to m(H)$. In the K\"unneth decomposition of $(f\times f)_*\Delta_*z$, $z\in H_n(L;k)$, all mixed positive-degree terms vanish. The remaining terms are $f_*z\otimes1+1\otimes f_*z$. Applying $q_*$ and conjugacy invariance gives $(jf)_*z=2(jf)_*z$, hence $(jf)_*z=0$. Induction shows that $n$ consecutive mitoses kill field homology through degree $n$. Every bar cycle in the union occurs at a finite stage, so all positive field homology of $M(H)$ vanishes. Rational vanishing makes its integral positive homology torsion. The universal-coefficient sequence with $\mathbb F_p$ in the next degree says that this torsion group has no elements of order $p$, for any prime $p$. It is zero.

This is the mitosis proof checked against \citet[Proposition 4.1, Theorem 4.2, and Lemma 5.4]{RootBDH1980}. For any group an equation $h^2=h$ implies $h=1$ by multiplication with $h^{-1}$, proving the final statement. No simplicity assertion is needed for the counterexamples; the stronger arbitrary-cardinal simple-overgroup assertion in CATHOT is not used here.
\end{proof}

\begin{claim}{B.11}{Proved}{The one-object test required after CATHOT}
Let $G$ be any nontrivial integrally acyclic small group, and let $\mathsf BG$ denote its one-object category. For every fundamental localizer satisfying Loc (6 b), $\mathsf BG\to e$ belongs to $\W$, although it does not induce an isomorphism on fundamental groups. Thus such a localizer cannot require $\pi_1$-isomorphisms even on one-object sources.
\end{claim}
\begin{proof}
The classifying nerve realizes to the usual $K(G,1)$: the nerve of the translation category with object set $G$ and a unique arrow between each pair is contractible, and its free $G$-quotient is $N\mathsf BG$. Its fundamental group is therefore $G$. The two-cone category on $\mathsf BG$ has objects $*,u,t$, endomorphisms $G$ at $*$, unique arrows $*\to u$ and $*\to t$, identities at $u,t$, and no other arrows. The composites of any group element with either crossing arrow are that same crossing arrow. Its realization is $\Sigma BG$, which is simply connected and acyclic, hence contractible. Apply \claimref{B.5} to $\mathsf BG\to e$, or apply \claimref{B.6} directly. The same cover by the two cones gives the infinite-arrow, three-object version of \claimref{B.8}. All its comma categories at the new terminal points are explicitly the indicated group category or a terminal cone; no idempotent issue is involved.
\end{proof}

\begin{claim}{B.12}{Proved}{The amalgam is a homological suspension model}
Let $H\hookrightarrow G$ be an injective homomorphism of small groups with $G$ integrally acyclic. The double mapping cylinder of $BG\leftarrow BH\to BG$ is a classifying space for $D=G*_HG$. Collapsing its two vertex spaces gives a specified map $BD\to\Sigma BH$ inducing an integral homology isomorphism. In particular this holds when $H$ is acyclic, but it does not require that assumption, and it does not assert a weak homotopy equivalence.
\end{claim}
\begin{proof}
The amalgam normal form embeds both vertex groups and identifies the fundamental group of this graph of spaces with $D$. The universal cover is a tree of contractible vertex and edge spaces over the Bass--Serre tree. Attaching a terminal vertex of a finite subtree preserves contractibility, so each finite subtree union is contractible. A cellular sphere and a cellular nullhomotopy problem meet finitely many cells and therefore a finite subtree; the entire cover is weakly contractible. This proves the asserted classifying-space model. One may use mapping cylinders of the classifying-space maps so that every edge-space attachment is a cofibration.

The quotient map to $\Sigma BH$ collapses both copies of $BG$ to the two suspension points and is the identity on the intervening cylinder of $BH$. The natural Mayer--Vietoris sequences for these two homotopy pushouts compare by the identities on $BH$ and by $BG\to *$ on both vertex spaces. These vertex maps are integral homology isomorphisms, including degree zero. The exact-sequence comparison proves that the quotient induces isomorphisms in every degree. This verifies the assertion from CATHOT B.8 and identifies the precise map. For acyclic nontrivial $H$, the suspension is contractible while $D$ contains $G$; this illustrates why only homology equivalence is claimed.
\end{proof}

\begin{claim}{B.13}{Refuted}{The upper interval is not the reversible class}
The assertion that $\WZ\subseteq\W$ suffices for Loc (6 b) is false. In particular $\W_0$ contains $\WZ$ and fails the cofiber criterion.
\end{claim}
\begin{proof}
Let $g$ be the fold from the discrete two-object category $\{a_0,a_1\}$ to $e$. It is not in $\W_0$. The category $K(g)$ has four objects $a_0,a_1,b,t$, its four nonidentity arrows are $a_i\to b$ and $a_i\to t$, and composition only involves identities. Its realization is the four-edge circle, hence is nonempty and connected and therefore $\W_0$-aspherical. The code checks these data and the Betti numbers $(1,1)$. The inclusion $\WZ\subseteq\W_0$ follows from the degree-zero argument in \claimref{B.6}. Thus a minimality theorem for the reverse axiom is not an upward-closure theorem for that axiom.
\end{proof}
